\question[7]
\lstinline|guessmyname2| (es.\ 10) offusca il segreto con XOR a chiave \lstinline|0x5A| sui byte
\lstinline|{0x1e,0x15,0x14,0x7d,0x0e,0x7a,0x0a,0x1b,0x14,0x13,0x19}|.
(a) Decodifica la stringa in chiaro. (b) Perché \lstinline|readelf|/\lstinline|strings| non bastano più? (c) Con quale comando gdb (i386, cdecl) leggi il 2\textsuperscript{o} argomento di \lstinline|strcmp|, e qual è l'alternativa con \lstinline|ltrace|?
\begin{solution}
(a) XOR byte a byte con \lstinline|0x5A| $\to$ \textbf{DON'T PANIC}
(\lstinline|0x1e^0x5A=0x44| `D', \lstinline|0x15^0x5A=0x4F| `O', \lstinline|0x14^0x5A=0x4E| `N', \lstinline|0x7d^0x5A=0x27| apostrofo, \lstinline|0x0e^0x5A=0x54| `T', \lstinline|0x7a^0x5A=0x20| spazio, \lstinline|0x0a^0x5A=0x50| `P', \lstinline|0x1b^0x5A=0x41| `A', \lstinline|0x14^0x5A=0x4E| `N', \lstinline|0x13^0x5A=0x49| `I', \lstinline|0x19^0x5A=0x43| `C').\\
(b) In \lstinline|.rodata| ci sono solo i byte offuscati, non la stringa: l'analisi statica non rivela più nulla.\\
(c) \lstinline|break strcmp|, \lstinline|run|, poi \lstinline|x/s *(char**)($esp+8)|; in alternativa \lstinline|ltrace -e strcmp ./guessmyname2|, che stampa \lstinline|strcmp("...", "DON'T PANIC")|.
\end{solution}
